
\subsubsection{4.1 Research Questions}

\begin{itemize}
\item \textbf{RQ1:} Does task granularity (function-level vs. repository-level) produce a measurable difference in GRPO advantage variance for a 7B-class model?
\item \textbf{RQ2:} Is reward sparsity (near-zero positive reward rate) the mechanistic pathway for advantage collapse?
\item \textbf{RQ3:} Is any observed difference practically significant?
\end{itemize}

\subsubsection{4.2 Dataset Details}

\begin{table}[htbp]
\centering
\begin{tabular}{llll}
\toprule
Dataset & Hub Name & Split & Size \\
\midrule
APPS & codeparrot/apps & train & 5,000 \\
CodeContests & deepmind/code_contests & train & 13,328 \\
SWE-bench Verified & princeton-nlp/SWE-bench_Verified & test & 500 \\
\bottomrule
\end{tabular}
\end{table}

\subsubsection{4.3 Baseline Conditions (H-E1 Infrastructure)}

The curriculum GRPO infrastructure supports 4 conditions:

\begin{table}[htbp]
\centering
\begin{tabular}{ll}
\toprule
Condition & Description \\
\midrule
curriculum & Easy tiers (0–2) for steps 0–2500; hard tiers (3–4) for steps 2501–5000 \\
uniform & Random sampling across all tiers throughout training \\
easy_only & Only easy tiers (0–2) throughout training \\
hard_only & Only hard tiers (3–4) throughout training \\
\bottomrule
\end{tabular}
\end{table}

\subsubsection{4.4 Implementation Details}

The training framework uses TRL 1.3.0 GRPOTrainer. APPS and CodeContests require schema unification before use in a shared training loop. The H-M1 advantage variance experiment used CodeLlama-7b-Instruct-hf. The H-E1 smoke test used DeepSeek-Coder-7B-base. GPU isolation was enforced via CUDA_VISIBLE_DEVICES on a single NVIDIA H100 NVL.
